\documentclass[10pt,twocolumn]{article}
\usepackage[scaled]{helvet}
\renewcommand\familydefault{\sfdefault}
\usepackage[T1]{fontenc}
\usepackage[margin=0.7in]{geometry}
\usepackage{titlesec}
\usepackage{enumitem}
\usepackage{parskip}
\setlength{\columnsep}{0.24in}
\setlist[itemize]{leftmargin=1.2em, topsep=0.2em, itemsep=0.18em}
\titleformat{\section}{\large\bfseries}{\thesection}{1em}{}

\begin{document}
\title{\vspace{-1.6cm}GitHub Actions with AWS}
\author{AWS ML Study Notes}
\date{}
\maketitle
\small

\section*{Overview}
GitHub Actions is a CI and automation system that runs directly from a GitHub repository. When paired with AWS, it becomes a common control plane for testing, packaging, and deploying cloud workloads from version controlled changes.

For many teams, GitHub Actions is attractive because code review, merge policy, and build automation all live in the same developer workflow.

\section*{Core Concepts}
\begin{itemize}
\item Workflows are defined in YAML and contain jobs and steps that run on GitHub hosted or self hosted runners.
\item The most secure AWS integration pattern is usually OpenID Connect role assumption, which lets GitHub obtain short lived AWS credentials without storing static secrets.
\item Reusable workflows, environment protections, artifact passing, and matrix execution help scale the setup across multiple services or model variants.
\end{itemize}

\section*{How It Works}
\begin{itemize}
\item A pull request or push event triggers the workflow, which then runs tests, linters, packaging steps, and security checks.
\item If deployment is allowed, the workflow assumes an AWS role and performs actions such as pushing an image to ECR, updating ECS, invoking a SageMaker pipeline, or applying infrastructure changes.
\item Workflow status feeds back into the repository so release readiness is visible during code review rather than after a merge.
\end{itemize}

\section*{AI and ML Relevance}
\begin{itemize}
\item GitHub Actions is often used to build training and inference containers, validate ML API code, trigger retraining workflows, or gate model deployment on evaluation outputs.
\item It also helps when the ML platform spans more than AWS alone, because GitHub can coordinate multi system automation from one source change event.
\end{itemize}

\section*{Operational Notes}
\begin{itemize}
\item Pin third party actions to trusted versions, avoid overprivileged tokens, and prefer OIDC over long lived AWS keys.
\item Use environment protections and branch rules so production credentials are reachable only from reviewed code paths.
\item Do not overload one workflow with every responsibility. Smaller purpose built workflows are easier to secure and troubleshoot.
\end{itemize}

\section*{What To Remember}
\begin{itemize}
\item GitHub Actions owns repository centered automation; AWS owns the runtime resources.
\item OIDC role assumption is the modern security baseline for GitHub to AWS federation.
\item For ML teams, the tool is most useful when connected to clear quality gates rather than only to deployment mechanics.
\end{itemize}

\end{document}
